%---
% slug: pptrain
% title: pptrain
% category: product
% eyebrow: TASK-AGNOSTIC LEARNING
% summary: A reusable library for training language models on synthetic tasks before ordinary pretraining and evaluating what transfers downstream.
% draft: false
% date: 2026-06-18
% rank: 3
%---
\documentclass[11pt]{article}
\usepackage{publication-web}
\begin{document}

\section{pptrain: ready-to-use prepretraining}

\textbf{Axym Labs}\\
Darmstadt, Germany

\paragraph{Task-agnostic training.}
Prepretraining is an emerging form of training that precedes ordinary pretraining. We term the broader objective \emph{task-agnostic training}: learning general patterns from synthetic or procedural tasks that can transfer across diverse downstream settings. Cellular automata, symbolic transformations, and mathematical primitives all fit this pattern.

\textbf{\href{https://github.com/Axym-Labs/pptrain}{Code}} \quad
\textbf{\href{https://pypi.org/project/pptrain/}{Package}}

One hurdle toward using these ideas in practice is that most implementations live in paper-specific code, hidden data generators, and one-off training scripts. Another is verifying that the pre-pretrainer was instantiated correctly and has a positive effect. \href{https://github.com/Axym-Labs/pptrain}{\texttt{pptrain}} provides a common interface for both.

\section{Ready to use}

\begin{verbatim}
pip install pptrain
\end{verbatim}

\begin{verbatim}
from pptrain import PrePreTrainer, RunConfig, create_task
from pptrain.integrations import HFCausalLMAdapter, HFModelConfig

trainer = PrePreTrainer(
    task=create_task("simpler_tasks", {"preset": "paper_binary_1m"}),
    model_adapter=HFCausalLMAdapter(
        HFModelConfig(model_name_or_path="sshleifer/tiny-gpt2")
    ),
    run_config=RunConfig(output_dir="runs/initial", max_steps=20),
)

bundle = trainer.fit().load_transfer_bundle()
\end{verbatim}

\section{Task families}

The library packages six task families from the literature as reusable presets: cellular-automata rollouts, balanced-bracket languages, procedural programs, symbolic operations, mathematical-reasoning primitives, and synthetic summarization. Each preset can be customized, and new task families can implement the same interface.

\section{Reference parity and evaluation}

Reference-parity utilities compare local outputs against original-source fixtures, because a changed generator, tokenizer, or serialization no longer tests the mechanism from the paper. The library also supports downstream transfer evaluation, compute-matched natural warmups, diagnostic metrics, plots, CSVs, and reports.

\end{document}
